\section*{الفصل \ref{ch:b1:faults-folds-fabrics} --- الصدوع والطيّات والنسيج الصخري}

\begin{solution}{exo:b1:faults-folds-fabrics:1}
(أ) \omterm{def:b1:faults-folds-fabrics:fault-types}{صدع عادي}. (ب) \omterm{def:b1:faults-folds-fabrics:fault-types}{صدع معكوس}، شبه انزلاق شاقولي خالص. (ج) \omterm{def:b1:faults-folds-fabrics:fault-types}{صدع انزياحي} \omterm{def:b1:faults-folds-fabrics:fault-types}{أيسر}. (د) \omterm{def:b1:faults-folds-fabrics:fault-types}{صدع انزياحي} \omterm{def:b1:faults-folds-fabrics:fault-types}{أيمن}. (ه) \omterm{def:b1:faults-folds-fabrics:fault-types}{صدع مائل الانزياح}:
\omterm{def:b1:faults-folds-fabrics:fault-elements}{الجدار المعلّق} نزل وتحرّك على طول \omterm{def:b1:geological-maps:attitude}{الامتداد} نحو اليسار؛ ومركّبته الانزياحية، $\cos\ang{40} = 0.77$ من \omterm{def:b1:faults-folds-fabrics:fault-elements}{الانزلاق الصافي}، أكبر من مركّبته الشاقولية، $0.64$، فيُسمّى صدعًا \omterm{def:b1:faults-folds-fabrics:fault-types}{أيسر} ذا مركّبة عادية.
(و) عند انحدار قدره \ang{25}، وهو أقلّ من \ang{45}، يصبح \omterm{def:b1:faults-folds-fabrics:fault-types}{الصدع المعكوس} صدعًا دفعيًّا.
\end{solution}

\begin{solution}{exo:b1:faults-folds-fabrics:2}
مع $|\sin r| = 1$، \omterm{def:b1:cross-sections:throw}{الإزاحة الرأسية} هي $300\sin\ang{55} = \qty{245.7}{m}$ \omterm{def:b1:cross-sections:throw}{والإزاحة الأفقية}
$300\cos\ang{55} = \qty{172.1}{m}$.
\end{solution}

\begin{solution}{exo:b1:faults-folds-fabrics:3}
\omterm{def:b1:faults-folds-fabrics:fold-classes}{لطية قائمة} ينحدر جناحاها بالزاويتين $\delta_1$ و$\delta_2$ في اتجاهين متعاكسين، $\iota = \ang{180} - \delta_1 - \delta_2$. (أ) \ang{110}: مفتوحة.
(ب) \ang{35}: ضيّقة. (ج) \ang{7}: متساوية الميل.
\end{solution}

\begin{solution}{exo:b1:faults-folds-fabrics:4}
الصدوع العادية تنحدر $\ang{45} + \ang{35}/2 = \ang{62.5}$، والدفعية $\ang{45} - \ang{17.5} = \ang{27.5}$. والصدوع الانزياحية المترافقة تصنع كل منها \ang{27.5} مع $\sigma_1$، على جانبيه، فتتلاقى عند \ang{55}.
\end{solution}

\begin{solution}{exo:b1:faults-folds-fabrics:5}
(أ) مركّبة الانزياح هي $400\cos\ang{30} = \qty{346.4}{m}$ على طول \omterm{def:b1:geological-maps:attitude}{الامتداد}، ومركّبة الانزلاق الشاقولي $400\sin\ang{-30} = \qty{-200}{m}$، أي \qty{200}{m} نزولًا على \omterm{def:b1:geological-maps:attitude}{الانحدار}؛ \omterm{def:b1:cross-sections:throw}{والإزاحة الرأسية} هي $200\sin\ang{50} = \qty{153.2}{m}$ \omterm{def:b1:cross-sections:throw}{والإزاحة الأفقية} $200\cos\ang{50} = \qty{128.6}{m}$. (ب) \omterm{def:b1:geological-maps:attitude}{الامتداد} هو $000$، ف$\mathbf s = (0, 1, 0)$؛ و$\mathbf u = (-\cos\ang{50}, 0, \sin\ang{50}) =
(-0.643, 0, 0.766)$. ثم
$\overrightarrow{AA'} = 400(0.866\,\mathbf s - 0.5\,\mathbf u) =
(128.6, 346.4, -153.2)$ بالأمتار. واتجاهها $\arctan(128.6/346.4) =
\ang{20.4}$، وغوصها $\arcsin(153.2/400) = \ang{22.5}$، وهو أيضًا
$\arcsin(\sin\ang{30}\sin\ang{50})$. (ج) \omterm{def:b1:faults-folds-fabrics:fault-types}{صدع مائل الانزياح}: \omterm{def:b1:faults-folds-fabrics:fault-elements}{زاوية الانزلاق} بين $0$ و\ang{-90}، أقرب إلى $0$: صدع \omterm{def:b1:faults-folds-fabrics:fault-types}{أيسر} ذو مركّبة عادية.
\end{solution}

\begin{solution}{exo:b1:faults-folds-fabrics:6}
(أ) يصنع الصدعان \ang{60} مع بعضهما؛ و$\sigma_1$ ينصّف الزاوية الحادّة، فهو أفقي على طول $050$ (و$230$)؛ و$\sigma_3$ أفقي على طول $140$؛ و$\sigma_2$ شاقولي، على الخطّ الذي يتلاقى فيه الصدعان الشاقوليان. ويصنع كل صدع \ang{30} مع $\sigma_1$، ف$\ang{45} - \phi/2 =
\ang{30}$ و$\phi = \ang{30}$. (ب) يقسّم الصدعان الأرض إلى أربع إسافين؛ الإسافين التي تحتوي $\sigma_1$، نحو $050$ ونحو $230$، تتحرّك نحو بعضها، وتلك التي تحتوي $\sigma_3$ تتباعد، نحو $140$ و$320$. خذ الصدع الذي امتداده $020$، بين الإسفين نحو $050$ (إلى شرقه) والإسفين نحو $320$ (إلى غربه): بالنسبة للإسفين الغربي، يتحرك الإسفين الشرقي نحو $230$ ويبتعد عن $320$، وكلاهما له مركّبة نحو $200$ على طول الصدع. فمراقب على الإسفين الغربي، ناظرًا نحو الشرق، يرى الإسفين الشرقي يتحرك نحو $200$، أي نحو اليمين: فالصدع الذي امتداده $020$ \omterm{def:b1:faults-folds-fabrics:fault-types}{أيمن}. والحجة نفسها، أو تماثل الحول $\sigma_1$، يجعل الصدع الذي امتداده $080$ \omterm{def:b1:faults-folds-fabrics:fault-types}{أيسر}.
\end{solution}

\begin{solution}{exo:b1:faults-folds-fabrics:7}
الناظمتان العلويتان هما
\begin{align*}
  \mathbf n_1 &= (\sin\ang{40}\sin\ang{120}, \sin\ang{40}\cos\ang{120}, \cos\ang{40}) = (0.557, -0.321, 0.766),\\
  \mathbf n_2 &= (\sin\ang{30}\sin\ang{260}, \sin\ang{30}\cos\ang{260}, \cos\ang{30}) = (-0.492, -0.087, 0.866).
\end{align*}
جداؤهما الاتجاهي هو
$(-0.211, -0.859, -0.207)$، وهو يشير نحو الأسفل: يغوص \omterm{def:b1:faults-folds-fabrics:fold-elements}{خطّ المفصل} بزاوية قدرها
$\arcsin 0.207/0.909 = \ang{13.1}$ نحو $\ang{180} + \arctan(0.211/0.859) =
\ang{193.8}$. $\mathbf n_1\cdot\mathbf n_2 = 0.417$، والزاوية بين العموديين
هي \ang{65.3}، و$\iota = \ang{114.7}$: \omterm{def:b1:faults-folds-fabrics:fold-classes}{طية مفتوحة}.
\end{solution}

\begin{solution}{exo:b1:faults-folds-fabrics:8}
$\mathbf n_1 = (0.321, 0.383, 0.866)$ للتطبّق، و
$\mathbf n_2 = (0.337, -0.925, 0.174)$ للتشقّق. وجداؤهما الاتجاهي، مأخوذًا نحو الأسفل، يعطي خطًّا يغوص \ang{25.4} نحو \ang{074.8}. وبحسب
\cref{prop:b1:faults-folds-fabrics:intersection}، فهذا هو اتجاه \omterm{def:b1:faults-folds-fabrics:fold-elements}{خطّ المفصل}: \omterm{def:b1:geological-maps:fold}{والطية} تغوص نحو \ang{25} تقريبًا نحو الشرق-الشمال-الشرقي، وهو ما ينبغي أن تؤكّده أيّ \omterm{def:b1:geological-maps:geological-map}{انكشاف} آخر لأحد الجناحين.
\end{solution}

\begin{solution}{exo:b1:faults-folds-fabrics:9}
(أ) مراقب جنوب الصدع، ناظرًا نحو الشمال، يرى الجزء الشمالي من العِرق مزاحًا نحو الشرق، أي نحو اليمين: فالصدع \omterm{def:b1:faults-folds-fabrics:fault-types}{أيمن}. (ب) العِرق مستوٍ شاقولي عمودي على الصدع، فإزاحته على طول \omterm{def:b1:geological-maps:attitude}{الامتداد} هي المركّبة الانزياحية: \qty{350}{m}، وهي الانزلاق كله إن كانت \omterm{def:b1:faults-folds-fabrics:fault-rocks}{الخدوش المنزلقة} أفقية. (ج) والانزلاق يصنع حينئذٍ \ang{20} مع \omterm{def:b1:geological-maps:attitude}{الامتداد} في مستوى الصدع الشاقولي: ومركّبته الأفقية ما زالت \qty{350}{m}، فيكون \omterm{def:b1:faults-folds-fabrics:fault-elements}{الانزلاق الصافي} $350/\cos\ang{20} = \qty{372.5}{m}$ ومركّبته الشاقولية $350\tan\ang{20} = \qty{127.4}{m}$. فالكتلة الشمالية تحرّكت نحو الشرق، والخدوش تنزل نحو الشرق، فهي إذن نزلت أيضًا، بمقدار \qty{127.4}{m}. والعِرق الشاقولي لا يستطيع أن يُظهر هذه الحركة الشاقولية، لأن انزلاقه على طول نفسه يترك أثره دون تغيير.
\end{solution}

\begin{solution}{exo:b1:faults-folds-fabrics:10}
(أ) العمودي على الصدع يصنع $\ang{90} - \theta$ مع $\sigma_1$، فبحسب \cref{prop:b1:stress-strain-rheology:mohr}،
$\sigma_n' = \sigma_1'\sin^2\theta + \sigma_3'\cos^2\theta$ و
$\tau = (\sigma_1' - \sigma_3')\sin\theta\cos\theta$. ويبدأ الانزلاق حين
$\tau = \mu\sigma_n'$؛ بالقسمة على $\sigma_3'\cos^2\theta$ وكتابة
$R = \sigma_1'/\sigma_3'$، يكون $(R - 1)\tan\theta = \mu(R\tan^2\theta + 1)$، ف
$R\tan\theta(1 - \mu\tan\theta) = \tan\theta + \mu$، وهذه هي الصيغة.
(ب) مع $\sigma_1$ شاقوليًّا، فإن $\theta$ هي \ang{90} ناقص \omterm{def:b1:geological-maps:attitude}{الانحدار}:
$\theta = \ang{30}$ تعطي $R = 3.12$؛ و$\theta = \ang{45}$، $R = 4.00$؛
$\theta = \ang{55}$، $R = 9.92$. (ج) يتلاشى المقام حين
$\tan\theta = 1/\mu$، أي $\theta = \ang{59.0}$؛ وبعده تكون $R$ سالبة:
فلا نسبة لإجهادات فعّالة موجبة تستطيع أن تجعل الصدع ينزلق، مهما كبر
$\sigma_1'$. والصدوع العادية التي تنحدر بأقلّ من \ang{31.0} لا تستطيع أن تنزلق
تحت $\sigma_1$ شاقولي مع $\mu = 0.6$. (د) الصدع الأفضل توجّهًا يصنع
$\ang{45} - \phi/2$ مع $\sigma_1$. من أجل $\mu = 0.6$، $\phi = \ang{31.0}$:
الصدوع العادية تنحدر \ang{60.5}، والدفعية \ang{29.5}. ومن أجل $\mu = 0.85$،
$\phi = \ang{40.4}$: \ang{65.2} و\ang{24.8}.
\end{solution}

\begin{solution}{exo:b1:faults-folds-fabrics:11}
ناظرين نحو الشمال، الشرق إلى اليمين، والمستويات المنحدرة نحو الغرب ترتفع نحو اليمين. (أ) عند B ينحدر التطبّق بالجهة نفسها التي ينحدر بها \omterm{def:b1:faults-folds-fabrics:foliation}{التشقّق}، وبأشدّ منه: فبحسب \cref{prop:b1:faults-folds-fabrics:cleavage-bedding}(أ)، تقع B على جناح منقلب؛ وA وC، حيث التطبّق أقلّ شدّة، في وضعهما الصحيح. (ب) يصعد \omterm{def:b1:faults-folds-fabrics:foliation}{التشقّق} بالزاوية \ang{60}؛ وعند A وC يصعد التطبّق بالزاوية \ang{40}، والدوران الحادّ من أحدهما إلى الآخر مع عقارب الساعة، عبر \ang{20}، فيحرّك أعلى \omterm{def:b1:faults-folds-fabrics:foliation}{التشقّق} نحو الشرق؛ وعند B يصعد التطبّق بالزاوية \ang{80}، والدوران عكس عقارب الساعة، عبر \ang{20}، فيحرّكه نحو الغرب. (ج) يتغيّر الاتجاه بين A وB، ثم بين B وC: فهناك مفصل في كل فترة. وعند A، وهو في وضعه الصحيح، يتحرك أعلى \omterm{def:b1:faults-folds-fabrics:foliation}{التشقّق} نحو الشرق، نحو \omterm{def:b1:geological-maps:fold}{الطية} المحدّبة عند أعلى الجناح: \omterm{def:b1:faults-folds-fabrics:fold-elements}{فالمفصل} بين A وB \omterm{def:b1:geological-maps:fold}{طية} محدّبة، والذي بين B وC \omterm{def:b1:geological-maps:fold}{طية} مقعّرة. (د) بحسب برهان القضية، \omterm{def:b1:faults-folds-fabrics:fold-elements}{الأجنحة} الصحيحة الوضع تنحدر $\gamma - \iota/2 = \ang{40}$، \omterm{def:b1:faults-folds-fabrics:fold-elements}{والأجنحة} المنقلبة $\gamma + \iota/2 =
\ang{80}$، ف$\gamma = \ang{60}$، كما يقول \omterm{def:b1:faults-folds-fabrics:foliation}{التشقّق}، و$\iota =
\ang{40}$: \omterm{def:b1:geological-maps:fold}{طيّات} ضيّقة. والمستويات المحورية تمتدّ شمالًا وتنحدر \ang{60} نحو الغرب، وضعيتها $270/60$. \omterm{def:b1:faults-folds-fabrics:fold-elements}{والأجنحة} القصيرة تجاوزت الشاقولي وأعلاها نحو الشرق: \omterm{def:b1:geological-maps:fold}{فالطيات} تنقلب نحو الشرق، كما في شكل الفصل.
\end{solution}

\begin{solution}{exo:b1:faults-folds-fabrics:12}
(أ) في \omterm{def:b1:geological-maps:fold}{طية} شيفرونية متماثلة يصنع كل جناح الزاوية $\iota/2$ مع \omterm{def:b1:faults-folds-fabrics:fold-elements}{المستوى المحوري}، وهو عمودي على اتجاه القصر؛ فجناح طوله $l$ يمتدّ $l\sin(\iota/2)$ عبر \omterm{def:b1:geological-maps:fold}{الطيّات}، والطبقة كلها $L_0\sin(\iota/2)$. (ب) $\sin\ang{30} = 0.5$: قصر قدره \qty{50}{\percent}. (ج) $\sin(\iota/2) = 0.7$ يعطي $\iota = \ang{88.9}$.
(د) $\iota = \ang{156}$ و$\sin\ang{78} = 0.978$: قصر قدره \qty{2.2}{\percent}، وهو أيضًا $1 - \cos\ang{12}$.
\end{solution}

\begin{solution}{pb:b1:faults-folds-fabrics:1}
\textbf{1.} $\mathbf n = (\sin\delta\sin\alpha, \sin\delta\cos\alpha,
\cos\delta)$ يعطي
\[
  \mathbf n_1 = (0.20855, 0.03415, 0.97742), \qquad
  \mathbf n_2 = (-0.20855, 0.03415, 0.97742).
\]

\textbf{2.} $\mathbf n_1\cdot\mathbf n_2 = -0.04349 + 0.00117 + 0.95535 =
0.91302$، والزاوية بين العموديين \ang{24.1}.

\textbf{3.} بحسب \cref{prop:b1:faults-folds-fabrics:hinge}،
$\iota = \ang{180} - \ang{24.1} = \ang{155.9}$.

\textbf{4.} بين \ang{170} و\ang{90}: \omterm{def:b1:faults-folds-fabrics:fold-classes}{طية مفتوحة}.

\textbf{5.} ينحدر الجناحان نحو $080.7$ و$279.3$، لا نحو الشرق والغرب تمامًا، لأن \omterm{def:b1:geological-maps:fold}{الطية} تغوص: فأشدّ خطوطهما انحدارًا ليست في المقطع، الذي لا يُظهر إلا \omterm{def:b1:cross-sections:apparent-dip}{الانحدار الظاهري}. وفي مقطع شرق--غرب، بحسب \cref{thm:b1:cross-sections:apparent-dip}، $\tan\beta = \tan\ang{12.2}
\cos\ang{9.3} = 0.2134$، ف$\beta = \ang{12.0}$، و$\ang{180} -
2\times\ang{12.0} = \ang{155.9}$ (المقطع مائل ب\ang{2} عن ذلك المقطع، وهذا لا يغيّر شيئًا بهذه الدقة).

\textbf{6.} أجنحتها موازية \omterm{def:b1:faults-folds-fabrics:fold-elements}{لأجنحة} \omterm{def:b1:geological-maps:fold}{الطية} المحدّبة، فزاوية أجنحتها أيضًا \ang{155.9}.

\textbf{7.} $\mathbf n_1\times\mathbf n_2 = (0.03415\times0.97742 -
0.97742\times0.03415,\; 0.97742\times(-0.20855) - 0.20855\times0.97742,\;
0.20855\times0.03415 + 0.03415\times0.20855) = (0, -0.40768, 0.01424)$.
مأخوذًا نحو الأسفل، $(0, 0.40768, -0.01424)$: يمتد \omterm{def:b1:faults-folds-fabrics:fold-elements}{خطّ المفصل} نحو $000$ ويغوص
$\arctan(0.01424/0.40768) = \ang{2.0}$ نحو الشمال.

\textbf{8.} \omterm{def:b1:faults-folds-fabrics:fold-elements}{المستوى المحوري} يحتوي \omterm{def:b1:faults-folds-fabrics:fold-elements}{خطّ المفصل} ويُنصّف الجناحين؛
وعموده على طول $\mathbf n_1 - \mathbf n_2 = (0.41710, 0, 0)$، أفقي
وشرق--غرب. فهو شاقولي ويمتدّ شمال--جنوب: $090/90$. \omterm{def:b1:geological-maps:fold}{والطية} المحدّبة قائمة.

\textbf{9.} \omterm{def:b1:faults-folds-fabrics:fold-classes}{طية مفتوحة}، قائمة، تغوص \ang{2.0} نحو الشمال.

\textbf{10.} بحسب \cref{prop:b1:geological-maps:folds}، تقترب \omterm{def:b1:geological-maps:geological-map}{الانكشافات} \omterm{def:b1:geological-maps:fold}{لطية} محدّبة من بعضها في اتجاه \omterm{def:b1:field-lab-methods-1:trend-plunge}{الغوص}: نحو الشمال. ولبّها يحتوي أقدم الصخور: الحجر الرملي في قاعدة الغطاء، وحيث قطع النهر أعمق ما يكون، قاعدة الشست والغرانيت تحته.

\textbf{11.} \omterm{def:b1:faults-folds-fabrics:foliation}{التشقّق المحوري المستوي} شاقولي ويمتدّ شمال--جنوب، $090/90$. وبحسب \cref{prop:b1:faults-folds-fabrics:intersection}، فإن تخطّطه بالتقاطع مع التطبّق موازٍ \omterm{def:b1:faults-folds-fabrics:fold-elements}{لخطّ المفصل}: اتجاهه $000$، وغوصه \ang{2.0}. مباشرة: المستوى $080.7/12.2$ يلتقي المستوى الشاقولي الممتدّ شمالًا على طول خطّ غوصه $\arctan(\tan\ang{12.2}
\cos\ang{80.7}) = \ang{2.0}$ نحو $000$.

\textbf{12.} ناظرين نحو الشمال، يرتفع الجناح الشرقي نحو اليسار بزاوية \ang{12.0}، \omterm{def:b1:faults-folds-fabrics:foliation}{والتشقّق} شاقولي: فالدوران الحادّ من \omterm{def:b1:faults-folds-fabrics:foliation}{التشقّق} إلى التطبّق، عبر \ang{78.0}، عكس عقارب الساعة، ويحرّك أعلى \omterm{def:b1:faults-folds-fabrics:foliation}{التشقّق} نحو الغرب. والجناح في وضعه الصحيح، فبحسب \cref{prop:b1:faults-folds-fabrics:cleavage-bedding}(ب) يقع مفصل \omterm{def:b1:geological-maps:fold}{الطية} المحدّبة نحو الغرب، عند $x = \qty{3.5}{km}$، كما تبيّن الخريطة.

\textbf{13.} \omterm{def:b1:geological-maps:attitude}{الامتداد} هو $120 - 90 = 030$؛ والمستوى ينحدر نحو $120$، ف\,\,
\omterm{def:b1:faults-folds-fabrics:fault-elements}{الجدار المعلّق} يقع على الجانب شرق--جنوب--شرق من الأثر.

\textbf{14.} \omterm{def:b1:faults-folds-fabrics:fault-types}{الصدع عادي}: ف$\sigma_1$ كان شاقوليًّا.

\textbf{15.} $\sigma_2$ أفقي على طول \omterm{def:b1:geological-maps:attitude}{الامتداد}، $030$--$210$؛
و$\sigma_3$ أفقي عبره، $120$--$300$.

\textbf{16.} $\ang{45} + \phi/2 = \ang{60}$ يعطي $\phi = \ang{30}$، و
$\mu = \tan\ang{30} = 0.577$.

\textbf{17.} المجموعة المترافقة امتدادها $030$ أيضًا، وتنحدر \ang{60} في الاتجاه الآخر: $300/60$. والمجموعتان تتلاقيان عند $\ang{180} - 2\times\ang{60} = \ang{60}$، و$\sigma_1$ الشاقولي ينصّف الزاوية الحادّة.

\textbf{18.} مع $\mu = 0.85$، $\phi = \ang{40.4}$، وانحدار أفضل الصدوع العادية توجّهًا \ang{65.2}؛ وصدع فاردن، عند \ang{60}، يبعد عنه \ang{5.2}. وبحسب \cref{exo:b1:faults-folds-fabrics:10}، ينزلق عند النسبة
$\sigma_1'/\sigma_3' = (1 + 0.85\cot\ang{30})/(1 - 0.85\tan\ang{30}) = 4.85$،
مقابل $4.68$ للصدع الأفضل: فهو قريب جدًّا، لكن ليس الأفضل توجّهًا بالضبط، وكان سينزلق مجددًا عند نسبة إجهادات أعلى ب\qty{4}{\percent} فقط.

\textbf{19.} الانزلاق نحو أسفل \omterm{def:b1:geological-maps:attitude}{الانحدار}: $r = \ang{-90}$؛ واتجاه \omterm{def:b1:faults-folds-fabrics:fault-elements}{متجهة الانزلاق} $120$، وغوصها \ang{60}.

\textbf{20.} $(\cos\ang{60}\sin\ang{120}, \cos\ang{60}\cos\ang{120},
-\sin\ang{60}) = (0.4330, -0.2500, -0.8660)$.

\textbf{21.} $e = 200/\tan\ang{60} = \qty{115.5}{m}$.

\textbf{22.} شرقًا: $115.47\sin\ang{120} = \qty{100.0}{m}$؛ شمالًا:
$115.47\cos\ang{120} = \qty{-57.7}{m}$، أي \qty{57.7}{m} نحو الجنوب.

\textbf{23.} $L = 200/\sin\ang{60} = \qty{230.9}{m}$، و
$200^2 + 115.47^2 = 40\,000 + 13\,333 = 53\,333 = 230.94^2$. \omterm{def:b1:cross-sections:throw}{الإزاحة الرأسية} \omterm{def:b1:cross-sections:throw}{والإزاحة الأفقية} هما المركّبتان الشاقولية والأفقية للانزلاق، الواقع في المستوى الشاقولي للانحدار: وهما ضلعا مثلث قائم الزاوية وترُه هو الانزلاق.

\textbf{24.} لأقرب متر، \omterm{def:b1:faults-folds-fabrics:fault-elements}{الانزلاق الصافي} على صدع فاردن هو
\textbf{\qty[text-series-to-math=true]{231}{m}}.
\end{solution}
